% !TEX program = xelatex
% 2027 新领军学科能力二试 · 综合 · 经济（预测训练卷）
\documentclass[cjk=noto, seal=first]{mojing-exam}

\examtitle{2027 新领军学科能力测试 · 预测训练卷}
\secret{训练用卷 ★ 非官方试题}

\begin{document}

\begin{paper}[time={建议 1.5 小时（训练设置）}, full=100]{二试 · 综合 · 经济}

\begin{instructions}
  \item 官方只公布二试综合六选二的科目名称（经济与社科是两个独立科目），没有公布经济的考纲、题型与分值。本卷以高中《经济与社会》相关内容为基础，
        考查用简单模型分析经济问题的能力，所需概念均在题中给出，全部安排为训练设置。
  \item 本卷共 4 题，每题 25 分，满分 100 分。
\end{instructions}

\begin{gradingnotes}
  \item 本评分标准为训练用建议标准。分析题按要点给分；分值为累计得分。
\end{gradingnotes}

\scoretable{1,2,3,4}

\begin{questions}

% =============================================================================
\question[25]
（税收由谁负担）某商品的需求函数为 $Q_d=100-2P$，供给函数为 $Q_s=3P-20$（$P$ 为价格，$Q$ 为数量）。
\begin{enumerate}
  \item 求市场均衡价格与数量；
  \item 政府对卖方每售出一件征税 5 元。求新的均衡数量、买方支付的价格、卖方实际得到的价格和政府税收，以及买卖双方每件各承担了多少税负；
  \item 消费者剩余指买方愿意支付的最高价格与实际支付价格之差的总和（需求曲线以下、价格线以上的面积），生产者剩余指卖方实际得到的价格与其愿意接受的最低价格之差的总和
        （价格线以下、供给曲线以上的面积）。征税后两者减少的部分中，未转化为税收的部分称为“无谓损失”。求无谓损失；
  \item 若改为对买方每买一件征税 5 元，结果有何不同？由此说明：税负在买卖双方之间如何分配取决于什么，而不取决于什么。
\end{enumerate}
\answerspace{2cm}
\begin{solution}
(1)~$100-2P=3P-20$，$P^*=24$，$Q^*=52$。\score{5}

(2)~卖方得到 $P-5$，供给变为 $Q=3(P-5)-20$。$100-2P=3P-35$，买方价格 $P_b=27$，卖方得到 $22$，数量 $46$，税收 $5\times46=230$。
买方每件多付 3 元，卖方每件少得 2 元，税负 $3:2$。\score{14}

(3)~无谓损失为三角形面积 $\frac12\times5\times(52-46)=15$。\score{18}

(4)~结果完全相同：买方总支付 27，卖方得到 22，数量 46。税负分配取决于供求曲线的斜率（弹性）——需求斜率 $2$、供给斜率 $3$，
较缺乏弹性（对价格反应较小）的一方承担更多；与法律上由谁缴税无关。\score{25}
\end{solution}

% =============================================================================
\question[25]
（囚徒困境与重复博弈）两家企业同时选择“合作”（维持高价）或“背叛”（降价抢市场）。每期收益（行为企业甲，列为企业乙，括号内先甲后乙）为：
\begin{center}
\begin{tabular}{c|cc}
  \toprule
  & 合作 & 背叛 \\
  \midrule
  合作 & $(3,3)$ & $(0,5)$ \\
  背叛 & $(5,0)$ & $(1,1)$ \\
  \bottomrule
\end{tabular}
\end{center}
“纳什均衡”指：在对方策略给定时，任何一方单独改变策略都不能使自己收益增加。
\begin{enumerate}
  \item 求一次性博弈的纳什均衡，并说明它为什么对双方都不是最好的结果；
  \item 博弈无限次重复，双方对未来收益按每期贴现因子 $\delta\in(0,1)$ 计算总收益（第 $t$ 期收益乘 $\delta^{t}$）。考虑“冷酷策略”：
        先合作，一旦有任何一方背叛就从此永远背叛。求使“双方都采用冷酷策略”成为均衡（任一方单独偏离都无利可图）的 $\delta$ 的范围；
  \item 用 (2) 的结论解释：为什么同一市场中长期竞争的少数企业比较容易维持高价？
\end{enumerate}
\answerspace{2cm}
\begin{solution}
(1)~无论对方如何选择，背叛的收益都更高（$5>3$，$1>0$），唯一纳什均衡为（背叛，背叛），收益 $(1,1)$；而双方合作得 $(3,3)$，个体理性导致集体不理性。\score{9}

(2)~一直合作的总收益 $\frac{3}{1-\delta}$；在某期背叛，该期得 5，此后对方永远背叛，最好也只得每期 1：$5+\frac{\delta}{1-\delta}$。
不偏离当且仅当 $\frac3{1-\delta}\ge5+\frac{\delta}{1-\delta}$，即 $3\ge5(1-\delta)+\delta$，$\delta\ge\frac12$。\score{20}

(3)~少数企业长期共存，相当于重复博弈且对未来重视程度高（$\delta$ 大）。短期降价的好处会被今后长期的价格战抵消，于是“合作”可以自我维持；
这也说明为什么反垄断机构要防范默契合谋。\score{25}
\end{solution}

% =============================================================================
\question[25]
（贷款的两种还款方式）某人贷款 $P=100$ 万元，期限 30 年，按月计息，月利率 $i=0.5\%$，共还 $n=360$ 期。
\begin{enumerate}
  \item “等额本息”：每月还款额相同，记为 $A$。证明 $A=\dfrac{Pi}{1-(1+i)^{-n}}$，并计算 $A$ 与总利息（已知 $1.005^{360}\approx6.0226$）；
  \item “等额本金”：每月偿还本金 $P/n$，另付当月剩余本金的利息。求第一个月的还款额和总利息；
  \item 两种方式哪种总利息少？总利息少是否意味着它对借款人一定更有利？说明理由。
\end{enumerate}
\answerspace{2cm}
\begin{solution}
(1)~各期还款折算到借款时刻的现值之和等于本金：$P=\sum_{k=1}^nA(1+i)^{-k}=A\cdot\frac{1-(1+i)^{-n}}{i}$，即得公式。\score{7}
$A=\frac{10^6\times0.005}{1-1/6.0226}\approx5996$ 元，总利息 $\approx360\times5995.5-10^6\approx1.16\times10^6$ 元。\score{12}

(2)~第一个月 $\frac{10^6}{360}+10^6\times0.005\approx7778$ 元。第 $k$ 个月利息为 $\bigl(P-\frac{(k-1)P}n\bigr)i$，总利息
$Pi\cdot\frac{n+1}2=10^6\times0.005\times180.5=902\,500$ 元。\score{19}

(3)~等额本金总利息少，因为它还本更快、占用资金更少；但前期还款压力大。两种方式在月利率 $i$ 下的现值都恰好等于 $P$，按市场利率折现并无优劣；
哪种更有利取决于借款人的现金流、其他投资机会的收益率等。\score{25}
\end{solution}

% =============================================================================
\question[25]
（比较优势）甲、乙两国都只生产布和粮两种产品，劳动是唯一投入。每单位劳动的产量如下：
\begin{center}
\begin{tabular}{c|cc}
  \toprule
  & 布 & 粮 \\
  \midrule
  甲国 & 6 & 2 \\
  乙国 & 1 & 1 \\
  \bottomrule
\end{tabular}
\end{center}
\begin{enumerate}
  \item 一国生产某产品的单位劳动产量高于另一国，称它在该产品上有“绝对优势”；一国生产某产品的机会成本（为多生产 1 单位而必须放弃的另一产品数量）低于另一国，
        称它在该产品上有“比较优势”。哪国在两种产品上都有绝对优势？分别求两国生产 1 单位粮的机会成本（以布计），并判断各国在哪种产品上有比较优势；
  \item 两国按比较优势分工后交换。粮与布的交换比例（1 单位粮换多少布）在什么范围内，两国都能从贸易中获益？
  \item 设甲国有 100 单位劳动、乙国有 60 单位劳动。若交换比例为 1 单位粮换 2 单位布，举例说明两国通过分工与贸易都能消费到原先无法自产的组合；
  \item 指出上述模型的两条重要假设，并说明现实中贸易可能使一国内部某些群体受损的原因。
\end{enumerate}
\answerspace{2cm}
\begin{solution}
(1)~甲国两种产品都更有效率，有绝对优势。1 单位粮的机会成本：甲国 3 单位布，乙国 1 单位布。乙国生产粮的机会成本低，有粮的比较优势；甲国有布的比较优势。\score{8}

(2)~交换比例在 1 与 3 之间（1 单位粮换 $1\sim3$ 单位布）时双方都获益。\score{13}

(3)~例：甲国全部生产布 600，乙国全部生产粮 60。乙国用 30 单位粮换甲国 60 单位布：甲国消费布 540、粮 30；
自产时甲国若要 30 单位粮需 15 单位劳动，只剩 85 单位劳动产布 510，现可得 540。乙国消费粮 30、布 60，自产时 30 单位粮需 30 单位劳动，
剩 30 单位劳动只能产布 30，现可得 60。两国都更好。\score{20}

(4)~假设：劳动可在行业间自由、无成本地转移；不考虑运输成本、规模经济与汇率等（亦可答：充分就业、两种产品在两国同质）。
现实中要素转移有成本，进口竞争行业的工人与企业可能失业或收入下降，贸易的总收益并不会自动补偿受损者，需要政策调节。\score{25}
\end{solution}

\end{questions}
\end{paper}
\end{document}
